
\documentclass{article}
\usepackage{colortbl}
\usepackage{makecell}
\usepackage{multirow}
\usepackage{supertabular}

\begin{document}

\newcounter{utterance}

\centering \large Interaction Transcript for game `ta\_mastermind', experiment `Mastermind{-}v0', episode 5 with BSU{-}SLIM\_\_Qwen3.5{-}9B{-}Base\_\_playpen{-}prm{-}9b{-}best\_of\_n.
\vspace{24pt}

{ \footnotesize  \setcounter{utterance}{1}
\setlength{\tabcolsep}{0pt}
\begin{supertabular}{c@{$\;$}|p{.15\linewidth}@{}p{.15\linewidth}p{.15\linewidth}p{.15\linewidth}p{.15\linewidth}p{.15\linewidth}}
    \# & \multicolumn{2}{c}{Player} && \multicolumn{2}{c}{Game Master} \\
    \hline

    \theutterance \stepcounter{utterance}  
    & & & \multicolumn{4}{p{0.6\linewidth}}{
        \cellcolor[rgb]{0.9,0.9,0.9}{
            \makecell[{{p{\linewidth}}}]{
                \texttt{\tiny{[P1$\langle$GM]}}
                \texttt{You are a competitive game player. Make sure you read the game instructions carefully, and always follow the required format.} \\
\\ 
\texttt{{[}GAME{]} You are playing Mastermind.} \\
\texttt{You need to find the code that is 4 digits long, each digit from 1 to 6, with no duplicates.} \\
\texttt{In your response, you can submit your guess in the following format: '{[}2 1 4 5{]}'.} \\
\texttt{After each guess, you will receive feedback in the form of black and white pegs.} \\
\texttt{A black peg indicates a correct digit in the correct position, while a white peg indicates a correct digit in the wrong position.} \\
\texttt{You have 20 turns to guess the code.} \\
\\ 
            }
        }
    }
     \\ \\

    \theutterance \stepcounter{utterance}  
    & \multicolumn{4}{p{0.6\linewidth}}{
        \cellcolor[rgb]{0.561,0.737,0.561}{
            \makecell[{{p{\linewidth}}}]{
                \texttt{\tiny{[P1$\rangle$GM]}}
                \texttt{I accept the challenge. I have carefully reviewed the rules:} \\
\texttt{1.  **Code Length**: 4 digits.} \\
\texttt{2.  **Digits**: 1 to 6.} \\
\texttt{3.  **Constraint**: No duplicates (all 4 digits must be unique).} \\
\texttt{4.  **Feedback**: Black peg = correct digit in correct position; White peg = correct digit in wrong position.} \\
\texttt{5.  **Format**: Guess must be submitted as `'{[}A B C D{]}'`.} \\
\texttt{6.  **Limit**: 20 turns.} \\
\\ 
\texttt{I am ready to begin. Please provide the initial feedback for my first guess or confirm that we can start.} \\
\\ 
\texttt{**My first guess is:**} \\
\texttt{`{[}1 2 3 4{]}`} \\
            }
        }
    }
    & & \\ \\

    \theutterance \stepcounter{utterance}  
    & & & \multicolumn{4}{p{0.6\linewidth}}{
        \cellcolor[rgb]{0.9,0.9,0.9}{
            \makecell[{{p{\linewidth}}}]{
                \texttt{\tiny{[P1$\langle$GM]}}
                \texttt{{[}GAME{]} Submitted {[}1 2 3 4{]}. Feedback: 1 black peg(s), 2 white peg(s).} \\
            }
        }
    }
     \\ \\

    \theutterance \stepcounter{utterance}  
    & \multicolumn{4}{p{0.6\linewidth}}{
        \cellcolor[rgb]{0.561,0.737,0.561}{
            \makecell[{{p{\linewidth}}}]{
                \texttt{\tiny{[P1$\rangle$GM]}}
                \texttt{**Analysis of Feedback:**} \\
\texttt{My guess was `{[}1 2 3 4{]}`.} \\
\texttt{The feedback is **1 black peg, 2 white pegs**.} \\
\\ 
\texttt{*   **Black Pegs (1):** One digit is in the exact correct position. Since there are no duplicates in the code, only one digit matches both the value and the position.} \\
\texttt{*   **White Pegs (2):** Two digits are present in the code but are in the wrong positions.} \\
\texttt{*   **Total Correct Digits:** 1 (black) + 2 (white) = 3 digits.} \\
\texttt{*   **Digits NOT in the Code:** Since there are 4 positions and only 3 digits are accounted for (1 black + 2 white), **one digit from my guess is completely incorrect** (it is not in the code at all).} \\
\\ 
\texttt{**Deduction:**} \\
\texttt{*   My guess: `1 2 3 4`} \\
\texttt{*   I have 1 Black (correct position) and 2 White (correct value, wrong position).} \\
\texttt{*   Therefore, one of these four digits is a "false alarm" and does not exist in the secret code.} \\
\texttt{*   To maximize efficiency, I should eliminate the digit that yields the least amount of information or try to test a specific position against a known digit.} \\
\\ 
\texttt{However, since I don't know *which* digit is the false alarm yet, I need a strategy to isolate it.} \\
\texttt{Let's try to keep the position of the **Black** peg fixed (to confirm it's truly correct and not an illusion caused by overlapping) and change the others.} \\
\\ 
\texttt{Actually, a more robust strategy for the first move is often to guess a pattern that covers different positions. But given the feedback 1B 2W, I have specific info.} \\
\texttt{Let's try to keep the digit that gave the Black peg in that position and move the others around to see if the White pegs become Black or disappear.} \\
\\ 
\texttt{Wait, if I have 1B at position 1, and 2W elsewhere.} \\
\texttt{If I swap the non{-}black digits, the Black should stay Black (unless I swapped the actual correct digit with a wrong one, which I won't do).} \\
\texttt{If I move the digit that was giving a White peg to a new position:} \\
\texttt{*   If that digit was in the code (it must be, for a White peg), and I move it to a new position:} \\
\texttt{    *   If the new position is the spot of the Black peg {-}> The Black becomes a White (because the original black digit is now displaced) and the moved digit might become the new Black. The total peg count (1B 2W) would remain the same or change structure.} \\
\texttt{    *   If the new position is empty (or occupied by a digit not in the code) {-}> The White peg disappears for that digit.} \\
\\ 
\texttt{Let's try a guess that keeps the first digit (`1`) in place but rearranges the rest (`2 3 4`).} \\
\texttt{Actually, let's try to test if `1` is the Black peg or if `2`, `3`, or `4` is the Black peg.} \\
\texttt{It's safer to keep the Black peg fixed to ensure it's real, then test the Whites.} \\
\\ 
\texttt{**Strategy:**} \\
\texttt{1.  Assume the Black peg at position 1 (`1`) is correct.} \\
\texttt{2.  Take the digits that provided the White pegs (`2`, `3`, `4` in some order).} \\
\texttt{3.  Construct a new guess where `1` stays at position 1, and `2`, `3`, `4` are permuted in positions 2, 3, 4.} \\
\texttt{4.  If the feedback for the new guess is still 1B 2W, then `1` is definitely the Black peg, and none of the other three are in the wrong positions? No, that would mean 3B.} \\
\texttt{5.  If the feedback changes, we learn something.} \\
\\ 
\texttt{Let's try `{[}1 4 3 2{]}` (Permuting the non{-}black digits).} \\
\texttt{If `1` is fixed and `2 3 4` are permuted:} \\
\texttt{*   If `2` was a White (meaning `2` is in the code but not at pos 1):} \\
\texttt{    *   In my next guess `{[}1 4 3 2{]}`, `2` is now at position 4.} \\
\texttt{    *   If the code is `1 2 3 4`, then `2` is at pos 2. Moving it to pos 4 makes it a Black peg.} \\
\texttt{    *   So `1` (Black) stays Black. `2` becomes Black. `3` and `4` become White (since they are in wrong positions).} \\
\texttt{    *   Result: 2B 1W.} \\
\texttt{*   If `2` was} \\
            }
        }
    }
    & & \\ \\

    \theutterance \stepcounter{utterance}  
    & & & \multicolumn{4}{p{0.6\linewidth}}{
        \cellcolor[rgb]{0.9,0.9,0.9}{
            \makecell[{{p{\linewidth}}}]{
                \texttt{\tiny{[P1$\langle$GM]}}
                \texttt{{[}GAME{]} Submitted {[}1 4 3 2{]}. Feedback: 1 black peg(s), 2 white peg(s).} \\
            }
        }
    }
     \\ \\

    \theutterance \stepcounter{utterance}  
    & \multicolumn{4}{p{0.6\linewidth}}{
        \cellcolor[rgb]{0.561,0.737,0.561}{
            \makecell[{{p{\linewidth}}}]{
                \texttt{\tiny{[P1$\rangle$GM]}}
                \texttt{**Analysis of Feedback:**} \\
\texttt{My previous guess was `{[}1 2 3 4{]}`, which returned **1 Black, 2 White**.} \\
\texttt{My new guess is `{[}1 4 3 2{]}`. I kept the digit `1` in the first position (the position of the Black peg) but permuted the other three digits (`2`, `3`, `4`) into positions `2`, `3`, `4`.} \\
\texttt{The feedback remains **1 Black, 2 White**.} \\
\\ 
\texttt{*   **Logic:**} \\
\texttt{    *   Since the feedback count did not increase, the digit `1` is definitely the correct digit in the correct position (Position 1). If `1` were not the correct digit, moving the other digits would have caused the Black peg to disappear (if I had moved the actual Black digit) or changed its value. Keeping `1` fixed and seeing the Black count stay at 1 confirms `1` is correct at Position 1.} \\
\texttt{    *   Since the White count stayed at 2, the remaining three positions (`2`, `3`, `4`) contain two correct digits in the wrong positions.} \\
\texttt{    *   Crucially, because `1` is fixed and correct, and the total correct digits identified is still 3 (1 Black + 2 White), **two of the digits from the set \{2, 3, 4\} are in the code**, and **one digit from the set \{2, 3, 4\} is NOT in the code at all**.} \\
\texttt{    *   Specifically, in my guess `{[}1 4 3 2{]}`, the digit `2` is now in Position 4. Since the feedback is still 1B 2W, `2` must be one of the "White" digits (i.e., it is in the code but not at Position 4). This implies that in the secret code, `2` is located at either Position 2 or Position 3.} \\
\texttt{    *   Consequently, the digit that was *not* in the code is one of the other two positions (`3` or `4`) in the context of the permutation.} \\
\\ 
\texttt{**Next Step:**} \\
\texttt{I need to isolate which position holds the second and third correct digits.} \\
\texttt{I will swap `2` back to Position 2 but move `3` and `4` around differently to test if they can generate a Black peg in positions 2 or 3.} \\
\texttt{Let's try putting `4` in Position 2 and `3` in Position 3, keeping `2` out or moving it to a new spot?} \\
\texttt{Actually, let's try to put one of the suspected "White" digits (let's say `4`) into the position where the Black peg is? No, that would likely mess up the Black peg count unless `4` becomes the new Black.} \\
\texttt{Better strategy: Keep `1` at Position 1. Try to make one of the other positions turn into a Black peg to see if the total Black count increases to 2.} \\
\texttt{If I guess `{[}1 4 3 2{]}` and get 1B 2W, it means `2` is in the code (but not at pos 4, so it's at pos 2 or 3).} \\
\texttt{Let's try guessing `{[}1 4 2 3{]}`.} \\
\texttt{*   `1` is fixed (Black).} \\
\texttt{*   `4` is at Pos 2.} \\
\texttt{*   `2` is at Pos 3.} \\
\texttt{*   `3` is at Pos 4.} \\
\texttt{If `4` is in the code but not at Pos 2, it should yield a White peg (if it goes to Pos 4, which is empty or wrong).} \\
\texttt{If `2` is in the code and is actually at Pos 2 (its natural spot), it should yield a Black peg.} \\
\texttt{So if the code is `1 4 2 X`, guessing `1 4 2 3` might give 2B (if 4 is wrong at pos 2? No wait).} \\
\\ 
\texttt{Let's refine the deduction:} \\
\texttt{Code structure: `1` is at Pos 1.} \\
\texttt{Remaining slots: 2, 3, 4.} \\
\texttt{Digits available for remaining slots: Two from `\{2, 3, 4\}`, One is "false" (not in code).} \\
\texttt{My guess `{[}1 4 3 2{]}` gave 1B 2W.} \\
\texttt{{-} `1` is Black (Correct).} \\
\texttt{{-} `4` is White (Correct value, wrong position).} \\
\texttt{{-} `3` is White (Correct value, wrong position).} \\
\texttt{{-} `2` is... wait.} \\
\texttt{Total correct digits = 1 (Black) + 2 (White) = 3.} \\
\texttt{So `1`, `4`, and `3` are in the code. `2` is NOT} \\
            }
        }
    }
    & & \\ \\

    \theutterance \stepcounter{utterance}  
    & & & \multicolumn{4}{p{0.6\linewidth}}{
        \cellcolor[rgb]{0.9,0.9,0.9}{
            \makecell[{{p{\linewidth}}}]{
                \texttt{\tiny{[P1$\langle$GM]}}
                \texttt{{[}GAME{]} You attempted an invalid move. Reason: You have already guessed {[}1, 4, 3, 2{]}. Please try a different guess. Please resubmit a valid move and remember to follow the game rules to avoid penalties.} \\
            }
        }
    }
     \\ \\

    \theutterance \stepcounter{utterance}  
    & \multicolumn{4}{p{0.6\linewidth}}{
        \cellcolor[rgb]{0.561,0.737,0.561}{
            \makecell[{{p{\linewidth}}}]{
                \texttt{\tiny{[P1$\rangle$GM]}}
                \texttt{**Correction:**} \\
\texttt{I realize my previous manual deduction contained a logical error, and my attempt to guess `{[}1 4 2 3{]}` was flagged as a repeat of a previous pattern structure I had already tested (or confused the unique permutation requirement). The game system enforces that every guess must be unique and valid.} \\
\\ 
\texttt{Let's re{-}evaluate the state based on the feedbacks:} \\
\texttt{1.  **Guess `{[}1 2 3 4{]}`**: 1 Black, 2 White.} \\
\texttt{    *   This tells us there are exactly 3 correct digits in total.} \\
\texttt{    *   One digit is in the correct position (the Black peg).} \\
\texttt{    *   Two digits are correct values but in the wrong positions (the White pegs).} \\
\texttt{    *   One digit is completely absent.} \\
\\ 
\texttt{2.  **Guess `{[}1 4 3 2{]}`**: 1 Black, 2 White.} \\
\texttt{    *   I kept `1` at Position 1.} \\
\texttt{    *   I permuted `2`, `3`, `4` into positions `4`, `3`, `2`.} \\
\texttt{    *   Since the feedback (1B, 2W) remained identical, this confirms my deduction:} \\
\texttt{        *   `1` is definitely the Black peg (Position 1).} \\
\texttt{        *   The set of "White" digits (`2`, `3`, `4` in some order) contains exactly two correct values and one incorrect value.} \\
\texttt{        *   Specifically, since swapping `2` and `4` didn't change the count, the "incorrect" digit (the one not in the code) is likely one of the ones I moved.} \\
\texttt{    *   Actually, let's look closer at `{[}1 4 3 2{]}`.} \\
\texttt{        *   If `2` was the "incorrect" digit in `{[}1 2 3 4{]}`, then in `{[}1 4 3 2{]}`, `2` is now at Position 4. If `2` is not in the code, it yields 0 pegs. `4` is at Position 2. If `4` is in the code, it yields 1 White. `3` is at Position 3. If `3` is in the code, it yields 1 White. Total: 1B (`1`), 1W (`4`), 1W (`3`). Result: 1B 2W. This fits.} \\
\texttt{        *   If `4` was the "incorrect" digit in `{[}1 2 3 4{]}`, then in `{[}1 4 3 2{]}`, `4` is at Position 2. If `4` is not in the code, it yields 0 pegs. `2` is at Position 4. If `2` is in the code, it yields 1 White. `3` is at Position 3. If `3` is in the code, it yields 1 White. Total: 1B 2W. This also fits.} \\
\texttt{    *   So, I know `1` is fixed at Pos 1.} \\
\texttt{    *   I know two digits from `\{2, 3, 4\}` are in the code, and one is not.} \\
\texttt{    *   I don't yet know *which* one is the false alarm, nor do I know the positions of the two correct ones relative to each other.} \\
\\ 
\texttt{**Next Strategy:**} \\
\texttt{I need to change the permutation of the "White" digits to see if I can isolate the false alarm or find a second Black peg.} \\
\texttt{Since `{[}1 4 3 2{]}` was invalid (already tried), I must try a new permutation of `2`, `3`, `4`.} \\
\texttt{Possible permutations of `\{2, 3, 4\}`:} \\
\texttt{1.  `{[}2 3 4{]}` (Original {-} Invalid)} \\
\texttt{2.  `{[}1 4 3 2{]}` (Just tried {-} Invalid)} \\
\texttt{3.  `{[}1 4 2 3{]}` (Try this next?)} \\
\texttt{4.  `{[}1 3 4 2{]}`} \\
\texttt{5.  `{[}1 3 2 4{]}`} \\
\texttt{6.  `{[}1 2 4 3{]}`} \\
\\ 
\texttt{Let's try **`{[}1 3 4 2{]}`**.} \\
\texttt{*   Keep `1` at Pos 1 (Fixed Black).} \\
\texttt{*   Place `3` at Pos 2.} \\
\texttt{*   Place `4` at Pos 3.} \\
\texttt{*   Place `2` at Pos 4.} \\
\\ 
\texttt{**My guess is:**} \\
\texttt{`{[}1 3 4 2{]}`} \\
            }
        }
    }
    & & \\ \\

    \theutterance \stepcounter{utterance}  
    & & & \multicolumn{4}{p{0.6\linewidth}}{
        \cellcolor[rgb]{0.9,0.9,0.9}{
            \makecell[{{p{\linewidth}}}]{
                \texttt{\tiny{[P1$\langle$GM]}}
                \texttt{{[}GAME{]} Submitted {[}1 3 4 2{]}. Feedback: 1 black peg(s), 2 white peg(s).} \\
            }
        }
    }
     \\ \\

    \theutterance \stepcounter{utterance}  
    & \multicolumn{4}{p{0.6\linewidth}}{
        \cellcolor[rgb]{0.561,0.737,0.561}{
            \makecell[{{p{\linewidth}}}]{
                \texttt{\tiny{[P1$\rangle$GM]}}
                \texttt{**Analysis of Feedback:**} \\
\texttt{*   **Previous Guess:** `{[}1 2 3 4{]}` → **1 Black, 2 White**} \\
\texttt{*   **Current Guess:** `{[}1 3 4 2{]}` → **1 Black, 2 White**} \\
\\ 
\texttt{**Deduction:**} \\
\texttt{1.  **Confirmed Black Peg:** The digit `1` at Position 1 is definitely correct. The Black peg count has remained at 1 despite changing the other three digits. If `1` were not the correct digit, swapping the others might have caused the Black peg to disappear (if I had accidentally swapped the true black digit) or changed its location. Since the count is stable, `1` is locked in Position 1.} \\
\texttt{2.  **Identifying the False Alarm:** In the first guess `{[}1 2 3 4{]}`, we had 2 White pegs. The digits available for the White pegs were `2`, `3`, and `4`. One of these is a "false alarm" (not in the code), and two are in the code but in the wrong positions.} \\
\texttt{    *   In `{[}1 2 3 4{]}`, `2` was at Pos 2, `3` at Pos 3, `4` at Pos 4. Result: 1B, 2W.} \\
\texttt{    *   In `{[}1 3 4 2{]}`, `3` was at Pos 2, `4` at Pos 3, `2` at Pos 4. Result: 1B, 2W.} \\
\texttt{    *   Since the result is identical to the first guess, the "false alarm" must be the digit that is in the same position in both guesses, OR the false alarm is one of the digits that I moved but the feedback didn't change because it remained a White peg in its new position?} \\
\texttt{    *   Wait, let's look at the positions of the *potential* White candidates:} \\
\texttt{        *   Candidate A: `2` (was at Pos 2 in Guess 1, at Pos 4 in Guess 2).} \\
\texttt{        *   Candidate B: `3` (was at Pos 3 in Guess 1, at Pos 2 in Guess 2).} \\
\texttt{        *   Candidate C: `4` (was at Pos 4 in Guess 1, at Pos 3 in Guess 2).} \\
\texttt{    *   If `2` were the false alarm (not in code):} \\
\texttt{        *   Guess 1: `2` (Pos 2) {-}> 0 pegs. `3` (Pos 3) {-}> 1W. `4` (Pos 4) {-}> 1W. Total: 1B, 2W.} \\
\texttt{        *   Guess 2: `2` (Pos 4) {-}> 0 pegs. `3` (Pos 2) {-}> 1W. `4` (Pos 3) {-}> 1W. Total: 1B, 2W.} \\
\texttt{        *   Result matches!} \\
\texttt{    *   If `3` were the false alarm:} \\
\texttt{        *   Guess 1: `3` (Pos 3) {-}> 0 pegs. `2` (Pos 2) {-}> 1W. `4` (Pos 4) {-}> 1W. Total: 1B, 2W.} \\
\texttt{        *   Guess 2: `3` (Pos 2) {-}> 0 pegs. `2` (Pos 4) {-}> 1W. `4` (Pos 3) {-}> 1W. Total: 1B, 2W.} \\
\texttt{        *   Result matches!} \\
\texttt{    *   It seems swapping any of the three candidates results in the same feedback because the set of positions 2, 3, 4 is just being permuted. This test doesn't isolate *which* digit is the false alarm yet.} \\
\\ 
\texttt{**Next Strategy:**} \\
\texttt{I need to break the symmetry. I need a guess where the positions of the potential White candidates change relative to each other or relative to the fixed `1` in a way that produces a different feedback pattern (e.g., 2 Black pegs, or 0 White pegs).} \\
\texttt{Since `1` is fixed at Pos 1, I need to try placing one of the digits (`2`, `3`, or `4`) at Position 1? No, `1` is already there, and there are no duplicates allowed. So `1` cannot be anywhere else.} \\
\texttt{I must place one of the other digits (`2`, `3`, or `4`) at Position 1? **Impossible**, because `1` is already confirmed at Position 1. The rules say "no duplicates".} \\
\\ 
\texttt{So, Position 1 is `1`.} \\
\texttt{Positions 2, 3, 4 must} \\
            }
        }
    }
    & & \\ \\

    \theutterance \stepcounter{utterance}  
    & & & \multicolumn{4}{p{0.6\linewidth}}{
        \cellcolor[rgb]{0.9,0.9,0.9}{
            \makecell[{{p{\linewidth}}}]{
                \texttt{\tiny{[P1$\langle$GM]}}
                \texttt{{[}GAME{]} You attempted an invalid move. Reason: You have already guessed {[}1, 3, 4, 2{]}. Please try a different guess. Please resubmit a valid move and remember to follow the game rules to avoid penalties.} \\
            }
        }
    }
     \\ \\

    \theutterance \stepcounter{utterance}  
    & \multicolumn{4}{p{0.6\linewidth}}{
        \cellcolor[rgb]{0.561,0.737,0.561}{
            \makecell[{{p{\linewidth}}}]{
                \texttt{\tiny{[P1$\rangle$GM]}}
                \texttt{**Correction:**} \\
\texttt{I see the issue. I have been trying to permute the set `\{2, 3, 4\}` across positions 2, 3, and 4. Since the feedback for `{[}1 2 3 4{]}` and `{[}1 4 3 2{]}` was identical (1B 2W), and now `{[}1 3 4 2{]}` is also 1B 2W, I have exhausted all permutations of the set `\{2, 3, 4\}` in those specific positions.} \\
\\ 
\texttt{This strongly suggests that my assumption about which digit is the "false alarm" is correct, but I cannot distinguish *which* one of `\{2, 3, 4\}` is the false alarm just by swapping them among themselves. The feedback remains constant because the "White" status of the false alarm doesn't change (it yields 0 pegs regardless of position 2, 3, or 4), and the "White" status of the two correct digits remains 1 each regardless of their swap.} \\
\\ 
\texttt{**New Strategy:**} \\
\texttt{I need to test if one of these digits is actually in the code by placing it in the **Black Position (Pos 1)**. Since `1` is already confirmed in Pos 1, I cannot place `1` anywhere else.} \\
\texttt{However, I can try to introduce a **new digit** into the mix to see if it yields a Black or White peg, which would help me eliminate candidates from `\{2, 3, 4\}`.} \\
\\ 
\texttt{Wait, if I replace one of the "suspected" digits with a completely new digit (e.g., `5` or `6`), the feedback will likely change significantly.} \\
\texttt{{-} If I replace `2` with `5`:} \\
\texttt{    {-} Guess: `{[}1 3 4 5{]}`} \\
\texttt{    {-} If `2` was the false alarm (not in code), then `5` is definitely not in the code either (since there is only room for 3 digits besides `1`, and I already have `3` and `4`). So `5` yields 0 pegs. `3` and `4` yield 2 White pegs. Result: 1B 2W.} \\
\texttt{    {-} If `2` was a correct digit (in code), then `5` is definitely not in the code (since `2` is in the code, and no duplicates). So `5` yields 0 pegs. `3` and `4` yield 2 White pegs. Result: 1B 2W.} \\
\texttt{    {-} This still yields 1B 2W. This doesn't help distinguish `2` from `3` or `4` either.} \\
\\ 
\texttt{**Re{-}evaluating the Logic:**} \\
\texttt{Let's step back.} \\
\texttt{Code has 4 unique digits from \{1..6\}.} \\
\texttt{Guess 1: `1 2 3 4` {-}> 1B 2W.} \\
\texttt{This means 3 digits are correct, 1 is wrong.} \\
\texttt{The correct digits are a subset of size 3 from `\{1, 2, 3, 4\}`.} \\
\\ 
\texttt{Possibility A: The wrong digit is `1`. (Impossible, because `1` is the only one that can provide the Black peg if I keep it fixed? No.)} \\
\texttt{Let's reconsider the Black Peg.} \\
\texttt{If the wrong digit is `1`, then the correct digits are from `\{2, 3, 4\}`.} \\
\texttt{If I guess `1 2 3 4` and get 1B 2W, and the wrong digit is `1`:} \\
\texttt{{-} `1` is not in code. So `1` at Pos 1 yields 0 Black, 0 White.} \\
\texttt{{-} `2` is in code. Where is it?} \\
\texttt{{-} `3` is in code. Where is it?} \\
\texttt{{-} `4` is in code. Where is it?} \\
\texttt{If the code has digits `2, 3, 4` plus one other digit (say `5`).} \\
\texttt{Guess `1 2 3 4`:} \\
\texttt{{-} `1` is wrong (0B, 0W).} \\
\texttt{{-} `2` is correct but at Pos 2. If its true position is 3 or 4, it gives a White peg.} \\
\texttt{{-} `3` is correct but at Pos 3. If its true position is 2 or 4, it gives a White peg.} \\
\texttt{{-} `4` is correct but at Pos 4. If its true position is 2 or 3, it gives a White peg.} \\
\texttt{Total: X Black, Y White.} \\
\texttt{If `1` is wrong, Max Black pegs = 1 (from the other three).} \\
\texttt{If I swap `2` and `3` to get `1 3 2 4` (Wait, I did} \\
            }
        }
    }
    & & \\ \\

    \theutterance \stepcounter{utterance}  
    & & & \multicolumn{4}{p{0.6\linewidth}}{
        \cellcolor[rgb]{0.9,0.9,0.9}{
            \makecell[{{p{\linewidth}}}]{
                \texttt{\tiny{[P1$\langle$GM]}}
                \texttt{{[}GAME{]} Submitted {[}1 3 4 5{]}. Feedback: 1 black peg(s), 2 white peg(s).} \\
            }
        }
    }
     \\ \\

    \theutterance \stepcounter{utterance}  
    & \multicolumn{4}{p{0.6\linewidth}}{
        \cellcolor[rgb]{0.561,0.737,0.561}{
            \makecell[{{p{\linewidth}}}]{
                \texttt{\tiny{[P1$\rangle$GM]}}
                \texttt{**Analysis of Feedback:**} \\
\texttt{*   **Guess 1 `{[}1 2 3 4{]}`**: 1 Black, 2 White.} \\
\texttt{*   **Guess 2 `{[}1 4 3 2{]}`**: 1 Black, 2 White.} \\
\texttt{*   **Guess 3 `{[}1 3 4 2{]}`**: 1 Black, 2 White.} \\
\texttt{*   **Guess 4 `{[}1 3 4 5{]}`**: 1 Black, 2 White.} \\
\\ 
\texttt{**Deduction:**} \\
\texttt{Every single guess I have made so far has resulted in `1` being at Position 1, and the feedback remaining exactly `1B, 2W`.} \\
\texttt{This is a very strong signal.} \\
\texttt{1.  **`1` is definitely correct and in Position 1.** If `1` were wrong (not in the code), removing it or moving it would eventually cause the Black peg count to drop to 0 (or change), because the Black peg count depends on the specific digit in the specific slot. Since every permutation with `1` at Pos 1 yields 1B, `1` must be the correct digit for that slot.} \\
\texttt{2.  **The "White" digits are `\{2, 3, 4\}`.** In all these guesses, the set of digits in positions 2, 3, and 4 has been `\{2, 3, 4\}` (just permuted). The feedback is consistently 2 White pegs.} \\
\texttt{    *   This implies that **all three** digits `2`, `3`, and `4` are present in the code (occupying positions 2, 3, and 4), but they are all in the wrong positions relative to the guess.} \\
\texttt{    *   Wait, let's look at the math again.} \\
\texttt{    *   If the code is `1 2 3 4` and I guess `1 3 4 5`:} \\
\texttt{        *   I assumed the code contained `1`, `3`, `4` and some other digit `5`.} \\
\texttt{        *   If `1` is in Pos 1 (Black).} \\
\texttt{        *   Then `3` is in Pos 2. If `3` is in the code, it's at some position. If I put `3` in Pos 2 and it yields a White peg, it means `3` is in the code but not at Pos 2.} \\
\texttt{        *   `4` is in Pos 3. If `4` is in the code and yields White, it's in the code but not at Pos 3.} \\
\texttt{        *   `5` is in Pos 4.} \\
\texttt{        *   Total: 1B (`1`), 2W (`3`, `4`). This works if the code is `1 5 2 3` (where `3` is at pos 3, `4` is at pos 4... wait).} \\
\texttt{    *   Let's re{-}evaluate the "False Alarm" theory.} \\
\texttt{    *   **Hypothesis:** The digit `1` is the **Wrong** digit (not in the code).} \\
\texttt{        *   If `1` is not in the code, then the Black peg in `1 2 3 4` must come from `2`, `3`, or `4`.} \\
\texttt{        *   If `1` is not in the code, then in `1 3 4 5`, `1` is still not in the code. So it contributes 0 pegs.} \\
\texttt{        *   The feedback is still 1B 2W. This means among `3`, `4`, `5`, there is 1 Black and 2 White.} \\
\texttt{        *   But in the first guess `1 2 3 4`, we had 1 Black and 2 White.} \\
\texttt{        *   This implies the "Black" provided by the set `\{2, 3, 4\}` in the first guess is the same as the "Black" provided by the set `\{3, 4, 5\}` in the fourth guess.} \\
\texttt{        *   This suggests the code might be `2 3 4 5` (all 4 digits).} \\
\texttt{        *   If the code is `2 3 4 5`:} \\
\texttt{            *   Guess `1 2 3 4`: `1` is wrong (0). `2` is at Pos 2 (Black, if code is `2`... wait, if code is `2 3 4 5`, then `2` is at Pos 1).} \\
\texttt{            *   If Code = `2 3 4 5`:} \\
\texttt{                *   Guess `1 2 3 4`: `1` (No), `2` (Yes, at Pos 1 {-}> Black} \\
            }
        }
    }
    & & \\ \\

    \theutterance \stepcounter{utterance}  
    & & & \multicolumn{4}{p{0.6\linewidth}}{
        \cellcolor[rgb]{0.9,0.9,0.9}{
            \makecell[{{p{\linewidth}}}]{
                \texttt{\tiny{[P1$\langle$GM]}}
                \texttt{{[}GAME{]} You attempted an invalid move. Reason: You have already guessed {[}1, 3, 4, 5{]}. Please try a different guess. Please resubmit a valid move and remember to follow the game rules to avoid penalties.} \\
            }
        }
    }
     \\ \\

    \theutterance \stepcounter{utterance}  
    & \multicolumn{4}{p{0.6\linewidth}}{
        \cellcolor[rgb]{0.561,0.737,0.561}{
            \makecell[{{p{\linewidth}}}]{
                \texttt{\tiny{[P1$\rangle$GM]}}
                \texttt{**Correction:**} \\
\texttt{I realize my previous analysis was leading me down a rabbit hole of permutations that were already tested or logically impossible under the strict "unique digits" constraint without clear differentiation. The game system is flagging repeats because I haven't sufficiently narrowed down the set of digits.} \\
\\ 
\texttt{Let's restart the logic cleanly.} \\
\texttt{1.  **Goal:** Find a code of 4 unique digits (1{-}6).} \\
\texttt{2.  **Constraints:** No duplicates.} \\
\texttt{3.  **Feedback:** Black (correct pos), White (correct val, wrong pos).} \\
\texttt{4.  **History:**} \\
\texttt{    *   `{[}1 2 3 4{]}` {-}> 1B 2W} \\
\texttt{    *   `{[}1 4 3 2{]}` {-}> 1B 2W} \\
\texttt{    *   `{[}1 3 4 2{]}` {-}> 1B 2W} \\
\texttt{    *   `{[}1 3 4 5{]}` {-}> 1B 2W} \\
\\ 
\texttt{**Observation:**} \\
\texttt{Every guess I have made so far has kept `1` at Position 1. The feedback has been identical (1B 2W).} \\
\texttt{This implies that `1` is **definitely** the correct digit for Position 1.} \\
\texttt{*   If `1` were *not* in the code, and I kept it at Position 1 in every guess, the Black peg count would likely fluctuate or drop to 0 when I swapped it out, or the "White" count would change unpredictably depending on where the *actual* correct digit happened to be. The fact that the pattern is stable suggests `1` is fixed.} \\
\\ 
\texttt{**The Problem:**} \\
\texttt{I have 1 Black peg and 2 White pegs.} \\
\texttt{Total correct digits = 3.} \\
\texttt{One digit is wrong (not in the code).} \\
\texttt{The other three are correct but in the wrong positions.} \\
\texttt{I have tested all permutations of the subset `\{2, 3, 4\}` in positions 2, 3, 4.} \\
\texttt{*   `{[}1 2 3 4{]}` {-}> 1B 2W (Pos 2=2, 3=3, 4=4)} \\
\texttt{*   `{[}1 4 3 2{]}` {-}> 1B 2W (Pos 2=4, 3=3, 4=2)} \\
\texttt{*   `{[}1 3 4 2{]}` {-}> 1B 2W (Pos 2=3, 4=4, 3=2)} \\
\texttt{*   `{[}1 3 4 5{]}` {-}> 1B 2W (Pos 2=3, 4=4, 5=5) {-}{-} Wait, I replaced `2` with `5`.} \\
\\ 
\texttt{**Deduction Logic Refinement:**} \\
\texttt{If `1` is fixed at Pos 1, and I swap the other digits around, the feedback remains 1B 2W.} \\
\texttt{This means:} \\
\texttt{*   The "White" pegs are generated by digits `2`, `3`, and `4`.} \\
\texttt{*   In every permutation of `\{2, 3, 4\}` (excluding `1`), I get 2 White pegs.} \\
\texttt{*   This implies that **all three digits `2`, `3`, and `4` are in the code**.} \\
\texttt{    *   Why? If one of them (say `2`) was *not* in the code, then in the guess `{[}1 3 4 2{]}`, `2` is at Pos 4. If `2` isn't in the code, it gives 0 pegs. `3` is at Pos 2. If `3` is in the code, it gives 1 White (since it's not at its true position... wait, where is it?).} \\
\texttt{    *   Actually, if `2` is not in the code, and `3`, `4` are, then in `{[}1 3 4 2{]}`:} \\
\texttt{        *   `1`: Black (Fixed).} \\
\texttt{        *   `3`: White (Correct val, wrong pos).} \\
\texttt{        *   `4`: White (Correct val, wrong pos).} \\
\texttt{        *   `2`: 0 (Not in code).} \\
\texttt{        *   Total: 1B 2W.} \\
\texttt{    *   So swapping `2` out gave 1B 2W.} \\
\texttt{    *   Now let's try removing `3`. Guess `{[}1 2 4 5{]}`.} \\
\texttt{        *   `1`: Black.} \\
\texttt{        *   `2`: White?} \\
\texttt{        *   `4`: White?} \\
\texttt{        *   `5`: 0?} \\
\texttt{        *   Result: 1B 2W.} \\
\texttt{    *   It seems as long as I have 3 correct digits and 1 wrong digit, and I} \\
            }
        }
    }
    & & \\ \\

    \theutterance \stepcounter{utterance}  
    & & & \multicolumn{4}{p{0.6\linewidth}}{
        \cellcolor[rgb]{0.9,0.9,0.9}{
            \makecell[{{p{\linewidth}}}]{
                \texttt{\tiny{[P1$\langle$GM]}}
                \texttt{{[}GAME{]} Submitted {[}1 2 4 5{]}. Feedback: 1 black peg(s), 2 white peg(s).} \\
            }
        }
    }
     \\ \\

    \theutterance \stepcounter{utterance}  
    & \multicolumn{4}{p{0.6\linewidth}}{
        \cellcolor[rgb]{0.561,0.737,0.561}{
            \makecell[{{p{\linewidth}}}]{
                \texttt{\tiny{[P1$\rangle$GM]}}
                \texttt{**Analysis of Feedback:**} \\
\texttt{*   **Guess 1 `{[}1 2 3 4{]}`**: 1B, 2W} \\
\texttt{*   **Guess 2 `{[}1 4 3 2{]}`**: 1B, 2W} \\
\texttt{*   **Guess 3 `{[}1 3 4 2{]}`**: 1B, 2W} \\
\texttt{*   **Guess 4 `{[}1 3 4 5{]}`**: 1B, 2W} \\
\texttt{*   **Guess 5 `{[}1 2 4 5{]}`**: 1B, 2W} \\
\\ 
\texttt{**Deduction:**} \\
\texttt{I have now tested every possible permutation of the set `\{2, 3, 4\}` in the last three positions, keeping `1` fixed at Position 1.} \\
\texttt{*   In every single case, the feedback has remained exactly **1 Black, 2 White**.} \\
\texttt{*   This consistency proves two things:} \\
\texttt{    1.  **`1` is definitely correct** and is in **Position 1**. If `1` were incorrect (i.e., not in the code), the Black peg would have disappeared or changed as soon as I stopped placing `1` in Position 1 (which I haven't done, so this is consistent), but more importantly, if I replaced `1` with another digit later, the pattern would break. But here, the stability across all permutations of the *other* digits confirms that the "White" pegs are generated by the set `\{2, 3, 4\}` regardless of their arrangement.} \\
\texttt{    2.  Since the feedback never changes when I swap `2`, `3`, and `4` around, it implies that **all three digits `2`, `3`, and `4` are present in the code**.} \\
\texttt{        *   *Reasoning:* If one of these digits (e.g., `2`) was *not* in the code, then in the guess where `2` is placed in a specific position (say Pos 4), that position would yield 0 pegs. In the next guess where `2` is moved to Pos 3, that position would still yield 0 pegs (since `2` isn't in the code). However, the other digits (`3`, `4`) would also shift positions. If the "White" count depends on how many of the placed digits match the code but are in the wrong spot, swapping a "0 peg" digit for another "0 peg" digit (if the other was also wrong) would keep the count, but mixing a "0 peg" digit with "1 peg" digits changes the total count of White pegs relative to the positions.} \\
\texttt{        *   Actually, simpler logic: The feedback is **1B 2W**. Total correct digits = 3.} \\
\texttt{        *   The code must contain `1`.} \\
\texttt{        *   The code must contain exactly **two** digits from the set `\{2, 3, 4\}`.} \\
\texttt{        *   Wait, if I tested *all* permutations of `\{2, 3, 4\}` and got the same result, does that mean all three are in the code?} \\
\texttt{        *   Let's trace:} \\
\texttt{            *   Case A: Code has `2` and `3` (and not `4`). Guess `{[}1 2 3 4{]}`: `4` is wrong (0W). `2` is at Pos 2 (White). `3` is at Pos 3 (White). Result: 1B 2W.} \\
\texttt{            *   Case B: Code has `2` and `4` (and not `3`). Guess `{[}1 2 3 4{]}`: `3` is wrong (0W). `2` is at Pos 2 (White). `4` is at Pos 4 (White). Result: 1B 2W.} \\
\texttt{            *   Case C: Code has `3` and `4` (and not `2`). Guess `{[}1 2 3 4{]}`: `2` is wrong (0W). `3` is at Pos 3 (White). `4` is at Pos 4 (White). Result: 1B 2W.} \\
\texttt{        *   **CRITICAL REALIZATION:** The feedback **1B 2W** is ambiguous regarding *which* two digits are correct among `\{2, 3, 4\}`. Just by permuting `\{2, 3, 4\}`, I cannot distinguish between the subsets `\{2,3\}`, `\{2,4\}`, or `\{3,4\}`.} \\
\texttt{        *   I need to test a guess that **excludes** a specific pair to see if the feedback changes (e.g., drops to 1B 1W or 1B 0W). This will eliminate the wrong pair} \\
            }
        }
    }
    & & \\ \\

    \theutterance \stepcounter{utterance}  
    & & & \multicolumn{4}{p{0.6\linewidth}}{
        \cellcolor[rgb]{0.9,0.9,0.9}{
            \makecell[{{p{\linewidth}}}]{
                \texttt{\tiny{[P1$\langle$GM]}}
                \texttt{{[}GAME{]} You attempted an invalid move. Reason: You have already guessed {[}1, 2, 3, 4{]}. Please try a different guess. Please resubmit a valid move and remember to follow the game rules to avoid penalties.} \\
            }
        }
    }
     \\ \\

    \theutterance \stepcounter{utterance}  
    & \multicolumn{4}{p{0.6\linewidth}}{
        \cellcolor[rgb]{0.561,0.737,0.561}{
            \makecell[{{p{\linewidth}}}]{
                \texttt{\tiny{[P1$\rangle$GM]}}
                \texttt{**Analysis of Situation:**} \\
\texttt{I have already tested all permutations of the set `\{2, 3, 4\}` in the remaining slots while keeping `1` fixed at Position 1.} \\
\texttt{{-} `{[}1 2 3 4{]}` → 1B 2W} \\
\texttt{{-} `{[}1 4 3 2{]}` → 1B 2W} \\
\texttt{{-} `{[}1 3 4 2{]}` → 1B 2W} \\
\texttt{The system flagged my next attempt (`{[}1 2 3 4{]}`) as a repeat.} \\
\texttt{This confirms that simply permuting the set `\{2, 3, 4\}` does not help me distinguish *which* subset of two digits is correct (e.g., is it `\{2,3\}` or `\{2,4\}`?). The feedback pattern `1B 2W` is identical for all these permutations.} \\
\\ 
\texttt{**New Strategy:**} \\
\texttt{I need to break the symmetry by introducing a **new digit** that hasn't been tested in the "variable" slots yet, specifically to test hypotheses about which digit is the "false alarm" (the one not in the code).} \\
\\ 
\texttt{Hypotheses for the set of 3 correct digits (given `1` is fixed):} \\
\texttt{1.  **Hypothesis A:** The correct digits are `\{1, 2, 3\}`. (So `4` is the false alarm).} \\
\texttt{2.  **Hypothesis B:** The correct digits are `\{1, 2, 4\}`. (So `3` is the false alarm).} \\
\texttt{3.  **Hypothesis C:** The correct digits are `\{1, 3, 4\}`. (So `2` is the false alarm).} \\
\\ 
\texttt{To test **Hypothesis A** (`4` is the false alarm):} \\
\texttt{{-} I need to construct a guess using `1`, `2`, `3` and a new digit (say `5` or `6`) to fill the 4th spot.} \\
\texttt{{-} Since `4` is not in the code, placing `4` anywhere will yield 0 pegs (0B, 0W).} \\
\texttt{{-} I need to place `2` and `3` in the remaining slots such that I get **2 White pegs**.} \\
\texttt{{-} If `2` and `3` are both in the code, one of them providing a Black peg and the other a White peg would result in **2B 1W**.} \\
\texttt{{-} If `2` and `3` are both in the code but neither is in the correct position for the Black peg, I would get **0B 2W**.} \\
\texttt{{-} Wait, I already know `1` is the Black peg (Position 1).} \\
\texttt{{-} So if I test Hypothesis A (`1, 2, 3` + `5`):} \\
\texttt{    {-} Pos 1: `1` (Matches real `1` → Black).} \\
\texttt{    {-} Pos 2: `2` (If `2` is in code but not at Pos 2 → White).} \\
\texttt{    {-} Pos 3: `3` (If `3` is in code but not at Pos 3 → White).} \\
\texttt{    {-} Pos 4: `5` (If `4` is the false alarm, `5` is not in code → 0 pegs).} \\
\texttt{    {-} **Predicted Feedback:** 1B, 2W.} \\
\texttt{    {-} **Result:** This matches my existing history! This means Hypothesis A is *possible*, but it doesn't prove it's the *only* one.} \\
\\ 
\texttt{To test **Hypothesis B** (`3` is the false alarm):} \\
\texttt{{-} I need to construct a guess using `1`, `2`, `4` and a new digit (say `5` or `6`).} \\
\texttt{{-} Place `3` in one of the slots to see if it acts as the "0 peg" digit.} \\
\texttt{{-} Let's try **`{[}1 2 3 5{]}`**? No, I already tried permutations of 2,3,4. I need to replace one of them with something else.} \\
\texttt{{-} Let's try **`{[}1 2 5 3{]}`**?} \\
\texttt{    {-} Pos 1: `1` (Black).} \\
\texttt{    {-} Pos 2: `2` (White? If `2` is in code and not at Pos 2).} \\
\texttt{    {-} Pos 3: `5` (If `3` is the false alarm, then `5` is not in code? No, if `3` is false, `5` must be in code).} \\
\texttt{    {-} Pos 4: `3` (If `3` is false, this yields 0).} \\
\texttt{    {-} If `3` is the false alarm, it contributes 0. `5` contributes 1 White. `2` contributes 1} \\
            }
        }
    }
    & & \\ \\

    \theutterance \stepcounter{utterance}  
    & & & \multicolumn{4}{p{0.6\linewidth}}{
        \cellcolor[rgb]{0.9,0.9,0.9}{
            \makecell[{{p{\linewidth}}}]{
                \texttt{\tiny{[P1$\langle$GM]}}
                \texttt{{[}GAME{]} Submitted {[}1 2 5 3{]}. Feedback: 2 black peg(s), 2 white peg(s).} \\
            }
        }
    }
     \\ \\

    \theutterance \stepcounter{utterance}  
    & \multicolumn{4}{p{0.6\linewidth}}{
        \cellcolor[rgb]{0.561,0.737,0.561}{
            \makecell[{{p{\linewidth}}}]{
                \texttt{\tiny{[P1$\rangle$GM]}}
                \texttt{**Analysis of Feedback:**} \\
\texttt{*   **Previous Guesses:** All permutations of `{[}1, 2, 3, 4{]}` yielded **1 Black, 2 White**.} \\
\texttt{*   **Current Guess:** `{[}1 2 5 3{]}` yielded **2 Black, 2 White**.} \\
\\ 
\texttt{**Deduction:**} \\
\texttt{This is the breakthrough!} \\
\texttt{1.  **The Black Peg Count Increased:** The feedback changed from 1B to 2B.} \\
\texttt{2.  **The Digit `1` is NOT the correct digit for Position 1.**} \\
\texttt{    *   In all previous tests, I placed `1` at Position 1 and got 1 Black peg.} \\
\texttt{    *   In the current test, I placed `1` at Position 1 and got 2 Black pegs.} \\
\texttt{    *   This implies that the digit at Position 1 is *not* `1`. If `1` were the true black digit, adding more correct digits (like moving `3` or `5` into the Black position) would require the feedback to change differently, or it implies `1` is wrong.} \\
\texttt{    *   Wait, let's look at the placement.} \\
\texttt{        *   Guess `{[}1 2 5 3{]}`:} \\
\texttt{            *   Pos 1: `1`} \\
\texttt{            *   Pos 2: `2`} \\
\texttt{            *   Pos 3: `5`} \\
\texttt{            *   Pos 4: `3`} \\
\texttt{        *   Result: 2 Black, 2 White.} \\
\texttt{    *   Since I have 2 Black pegs, two of my guessed digits are in the exact correct positions.} \\
\texttt{    *   Possibilities for the 2 Black pegs:} \\
\texttt{        *   Could Pos 1 (`1`) and Pos 2 (`2`) be correct? {-}> Code starts `1 2 ...`. But if `1` and `2` are correct in those spots, why did my previous guesses (like `{[}1 2 3 4{]}`) only get 1 Black?} \\
\texttt{            *   If Code is `1 2 X Y`:} \\
\texttt{            *   Guess `{[}1 2 3 4{]}`: `1` matches (Black), `2` matches (Black). Total 2B.} \\
\texttt{            *   But my history says `{[}1 2 3 4{]}` gave **1B 2W**.} \\
\texttt{            *   Therefore, the code **cannot** be `1 2 ...`. The configuration `1` at Pos 1 and `2` at Pos 2 is incorrect because it would have produced 2B in the first guess.} \\
\\ 
\texttt{        *   Could Pos 1 (`1`) and Pos 3 (`5`) be correct? {-}> Code starts `1 ... 5 ...`.} \\
\texttt{            *   If Code is `1 X 5 Y`:} \\
\texttt{            *   Guess `{[}1 2 5 3{]}`: `1` matches (Black), `5` matches (Black). Total 2B.} \\
\texttt{            *   Check previous guess `{[}1 2 3 4{]}`: `1` matches (Black). `2` at Pos 2 (White), `5` at Pos 3 (White), `3` at Pos 4 (White). Total 1B 3W?} \\
\texttt{            *   But my history says `{[}1 2 3 4{]}` gave **1B 2W**.} \\
\texttt{            *   So this is inconsistent.} \\
\\ 
\texttt{        *   **Conclusion:** The pattern of 2 Blacks in `{[}1 2 5 3{]}` combined with the history of 1B in all `{[}1...{]}` guesses means that **`1` is NOT the correct digit for Position 1**.} \\
\texttt{    *   Wait, if `1` is not correct for Pos 1, then the Black peg in `{[}1 2 5 3{]}` must come from `2` and `5`.} \\
\texttt{    *   So, `2` is at Pos 2 (Correct) AND `5` is at Pos 3 (Correct).} \\
\texttt{    *   This means the code starts with `2` and `5`.} \\
\texttt{    *   So the code is `2 5 ...`.} \\
\texttt{    *   But wait, if the code is `2 5 ...`, then in Guess `{[}1 2 3 4{]}`:} \\
\texttt{        *   Pos 1: `1` (Wrong, so 0B).} \\
\texttt{        *   Pos 2: `2` (Correct, so 1B).} \\
\texttt{        *   Pos 3: `3` (Wrong? If `5` is the correct digit for Pos 3, then `3` is wrong).} \\
\texttt{        *   Pos 4: `4` (Wrong?).} \\
\texttt{        *   Total: 1B 0W?} \\
\texttt{    *   But my history says} \\
            }
        }
    }
    & & \\ \\

    \theutterance \stepcounter{utterance}  
    & & & \multicolumn{4}{p{0.6\linewidth}}{
        \cellcolor[rgb]{0.9,0.9,0.9}{
            \makecell[{{p{\linewidth}}}]{
                \texttt{\tiny{[P1$\langle$GM]}}
                \texttt{{[}GAME{]} You attempted an invalid move. Reason: You have already guessed {[}1, 2, 3, 4{]}. Please try a different guess. Please resubmit a valid move and remember to follow the game rules to avoid penalties.} \\
            }
        }
    }
     \\ \\

    \theutterance \stepcounter{utterance}  
    & \multicolumn{4}{p{0.6\linewidth}}{
        \cellcolor[rgb]{0.561,0.737,0.561}{
            \makecell[{{p{\linewidth}}}]{
                \texttt{\tiny{[P1$\rangle$GM]}}
                \texttt{**Analysis of Feedback:**} \\
\texttt{*   **History:**} \\
\texttt{    *   `{[}1 2 3 4{]}` → 1B 2W} \\
\texttt{    *   `{[}1 2 5 3{]}` → 2B 2W} \\
\texttt{*   **Observation:**} \\
\texttt{    *   In `{[}1 2 5 3{]}`, I have **2 Black** pegs. This means two of the positions I guessed are exactly correct.} \\
\texttt{    *   In all previous guesses starting with `1` (e.g., `{[}1 2 3 4{]}`), I only had **1 Black** peg.} \\
\texttt{    *   Since I kept `1` at Position 1 in all previous guesses and only got 1 Black then, but now I have `1` at Position 1 and got 2 Blacks, the *second* Black peg must come from a different position than Position 1.} \\
\texttt{    *   Specifically, in `{[}1 2 5 3{]}`, the digits are at Pos 1, 2, 3, 4.} \\
\texttt{    *   If `1` (Pos 1) is NOT the correct digit for that position (because if it were, it would have contributed to the Black count in all the `{[}1...{]}` guesses, and since the count increased, it implies the *location* of the Black peg changed or a new one appeared).} \\
\texttt{    *   Actually, simpler logic:} \\
\texttt{        *   If `1` were the correct digit for Pos 1, then in `{[}1 2 5 3{]}`, `1` would give a Black peg.} \\
\texttt{        *   In `{[}1 2 3 4{]}`, `1` gave a Black peg.} \\
\texttt{        *   If `1` is correct at Pos 1, then in `{[}1 2 5 3{]}`, we have Black from `1`. To get 2 Blacks total, one of `2`, `5`, or `3` must also be in its correct position.} \\
\texttt{        *   But in `{[}1 2 3 4{]}`, we had 1 Black (from `1`). We also had 2 Whites. This means `2`, `3`, `4` were *not* in their correct positions (otherwise we would have had 2B).} \\
\texttt{        *   So, if `1` is correct at Pos 1, then `2`, `3`, `4` are *not* correct at Pos 2, 3, 4 respectively.} \\
\texttt{        *   Now look at `{[}1 2 5 3{]}`. `1` is at Pos 1 (Black). `2` is at Pos 2. If `2` is *not* correct at Pos 2 (from previous deduction), it should give 0 or White.} \\
\texttt{        *   If `5` is at Pos 3. If `3` is at Pos 4.} \\
\texttt{        *   If the code has `1` at Pos 1, then `2`, `3`, `4` are the other digits.} \\
\texttt{        *   If `2` is not at its correct position (Pos 2), and `3` is not at its correct position (Pos 3), and `4` is not at its correct position (Pos 4)... then we only have 1 Black.} \\
\texttt{        *   BUT, in `{[}1 2 5 3{]}`, I got 2 Blacks.} \\
\texttt{        *   This implies that **`1` is NOT the correct digit for Position 1**.} \\
\texttt{        *   If `1` is not correct at Pos 1, then the 2 Black pegs in `{[}1 2 5 3{]}` must come from the other positions.} \\
\texttt{        *   The other positions are 2, 3, and 4.} \\
\texttt{        *   So, **Position 2 and Position 3 are correct**.} \\
\texttt{        *   Therefore: **Pos 1 is wrong**. **Pos 2 is `2`**. **Pos 3 is `5`**.} \\
\texttt{        *   Wait, if Pos 2 is `2` and Pos 3 is `5`, then the code starts `2 5 ...`.} \\
\texttt{        *   Let's check the White pegs.} \\
\texttt{        *   In `{[}1 2 5 3{]}`, we have 2 Blacks (`2` at 2, `5` at 3).} \\
\texttt{        *   Remaining digits: `1` (at Pos 1) and `3` (at Pos 4).} \\
\texttt{        *   Since `1` is not at its correct position (we deduced Pos 1 is wrong for `1`), `1` at Pos 1 gives 0 pegs.} \\
\texttt{        *   Since `3` is at Pos 4. If `3` is in the code} \\
            }
        }
    }
    & & \\ \\

    \theutterance \stepcounter{utterance}  
    & & & \multicolumn{2}{p{0.3\linewidth}}{
        \cellcolor[rgb]{0.9,0.9,0.9}{
            \makecell[{{p{\linewidth}}}]{
                \texttt{\tiny{[GM$|$GM]}}
                \texttt{role: Player 0} \\
\texttt{invalid\_move: True} \\
\texttt{turn\_count: 7} \\
\texttt{reason: Invalid Move: You have already guessed {[}1, 2, 5, 3{]}. Please try a different guess.} \\
\texttt{reward: 0.75} \\
            }
        }
    }
    & & \\ \\

\end{supertabular}
}

\end{document}
